\documentclass[12pt]{article}
% Préambule commun à tous les documents extraits de « La Revue CAPES »
\usepackage[T1]{fontenc}
\usepackage[utf8]{inputenc}
\usepackage{lmodern}
\usepackage{amsmath,amssymb,amsthm,amsfonts,mathrsfs}
\usepackage{setspace}
\usepackage{listings}
\usepackage{pgfplots}
\pgfplotsset{compat=1.18}
\usepackage{longtable}
\usepackage{adjustbox}
\usepackage{lettrine}
\usepackage{tkz-tab}
\usepackage{changepage}
\usepackage{pdflscape}
\usepackage{graphicx}
\usepackage{tikz}
\usepackage{xcolor}
\usepackage[a4paper,top=2cm,bottom=2.5cm,left=2.5cm,right=2cm]{geometry}
\usepackage{fancyhdr}
\usepackage{draftwatermark}
\SetWatermarkText{La RevueCapes}
\SetWatermarkLightness{0.9}
\SetWatermarkScale{2}
\usepackage{hyperref}
\hypersetup{colorlinks=true,linkcolor=blue!50!black,urlcolor=blue!50!black}

\definecolor{revue}{RGB}{20,60,120}

\pagestyle{fancy}
\fancyhf{}
\renewcommand{\headrulewidth}{0.4pt}
\setlength{\headheight}{15pt}

% \entete{Matière}{Titre}{Type de document}
\newcommand{\entete}[3]{%
  \fancyhead[L]{\small\textsc{La Revue CAPES} -- #1}%
  \fancyhead[R]{\small #2}%
  \fancyfoot[C]{\thepage}%
  \thispagestyle{plain}%
  \begin{center}
    {\color{revue}\rule{\textwidth}{1.2pt}}\\[4pt]
    {\small\textsc{Concours directs CAP-CEG / CAPES -- Burkina Faso}}\\[6pt]
    {\LARGE\bfseries\color{revue} #2}\\[6pt]
    {\large\itshape #1 -- #3}\\[2pt]
    {\color{revue}\rule{\textwidth}{1.2pt}}
  \end{center}
  \vspace{0.5cm}%
}

% Encadré pour les remarques ajoutées dans les corrigés
\newcommand{\remarque}[1]{\par\medskip\noindent\fcolorbox{revue}{revue!6}{\parbox{\dimexpr\linewidth-2\fboxsep-2\fboxrule}{\textbf{\color{revue}Remarque.} #1}}\par\medskip}


\begin{document}
\entete{Mathématiques}{Sujet type n°9}{Énoncé}
\begin{spacing}{1.5}

 
 
 \medskip\textbf{Exercice 1}

On considère dans $\mathbb{R}^2$ le domaine $D$ délimité par les courbes d'équations $x = y^2/4$ et $y = 2x$.
 
 1. Calculer l'aire de $D$.
 
 2. Calculer l'intégrale

$\int\int_D(y-x)dxdy$.

3. On note $C$ le bord de $D$ orienté de sorte à laisser $D$ sur sa gauche. Calculer l'intégrale

$\int_Cy^2dx+xydy$

a. directement,

b. puis par la formule de Green-Riemann.


\medskip\textbf{Exercice 2}

1. Trouver la solution de l'équation différentielle

$y^{\prime\prime}+2y^{\prime}+y=2e^{-x}$ vérifiant $y(0)=3$ et $y^{\prime}(0)=1$

2. Résoudre les équations différentielles 

$(E_1): y^{\prime\prime}+y^{\prime}=x+chx$

$(E_2): y^{\prime\prime}-3y^{\prime}-4y=xe^{\mid x\mid}$

$(E_3): y^{\prime\prime}+y^{\prime}-2y=x^2e^{-2x}$

$(E_4): 4xy^{\prime\prime}+2y^{\prime}-y=0$


\medskip\textbf{Exercice 3}

Les question de cet exercice sont indépendantes.

1. Soit $F(X)=\frac{N(X)}{D(X)}$ une fraction rationnelle non nulle écrite sous forme irréductible. Montrer que si $a$ est un pôle simple de $F(X)$ alors $D'(a)\neq 0$ et que le coefficient de $\frac{1}{X-a}$ dans la décomposition en éléments simple de $F(X)$ est $\frac{N(a)}{D'(a)}$.

2. Donner la négation de la phrase : \emph{"Chaque ville du Burkina possède au moins une agence bancaire"}.

3. Donner la contraposée de la phrase \emph{"La campagne agricole ne sera pas une réussite si la pluviométrie n'est pas bonne"}

4. Dites si les affirmations suivantes sont vrais ou fausses. Justifiez votre réponse.

a. Soit $A,B$ deux matrices d'ordre $n$ : $AB=0_n\Longrightarrow A=0_n \hspace{0.1cm} ou \hspace{0.1cm} B=0_n$

b. La négation de l'assertion  $a\leq b$ ou $\mid a\mid> c$ est $a>b$ ou $\mid a\mid\leq c$.

5. Soient $E,F$ deux ensembles et $f:E\longrightarrow F$ une application. Soient $A\subset E$ et $B\subset F$.

Montrer que $f(A\cap f^{-1}(B))=f(A)\cap B$

\medskip\textbf{Exercice 4}

$\mathbb{C}$ désigne le corps des nombres complexes. 

Soit $f$ la fonction de $\mathbb{C}$ dans $\mathbb{C}$ définie, pour $z\neq 2i$, par : 
$f(z) = \frac{z -1}{z -2i}$

$M, A $ et $ B$ sont les points d'affixes respectives $z, 1, 2i$. 

1) Donner les formes cartésienne et trigonométrique de $f(i)$. 

2) Résoudre l'équation $f(z) = 2i$ 

3) Déterminer l'ensemble $D$ des points $M$ tels que $\mid f(z) \mid = 2$.

\quad
\quad

\end{spacing}
\end{document}
